\section*{Bab \ref{ch:b2:fourier} --- Deret Fourier}

\begin{solution}{exo:b2:fourier:1}
Keganjilannya membunuh $a_n$. Untuk $n \geq 1$:
\[
b_n = \frac{2}{\pi}\int_0^{\pi} \sin nt\,\dd t
= \frac{2}{\pi}\cdot\frac{1 - (-1)^n}{n}
= \begin{cases} \frac{4}{\pi n} & n \text{ ganjil},\\ 0 & n
\text{ genap}. \end{cases}
\]
Jadi $S(f)(t) = \frac{4}{\pi}\sum_{k\geq0} \frac{\sin\bigl((2k+1)t
\bigr)}{2k+1}$. Dirichlet di $t = \frac\pi2$ (yakni titik kekontinuan,
bernilai $1$): sebab $\sin\bigl((2k+1)\frac\pi2\bigr) = (-1)^k$, sehingga memberi
\[
1 = \frac4\pi \sum_{k\geq0}\frac{(-1)^k}{2k+1}
\quad\Longrightarrow\quad
\sum_{k\geq0}\frac{(-1)^k}{2k+1} = \frac{\pi}{4}
\quad\text{(Leibniz)}.
\]
Di lompatan $t = 0$: deretnya berjumlah $0 = \frac{f(0^+) +
f(0^-)}{2}$, sebagaimana ditetapkan Dirichlet (sebab setiap sukunya lenyap:
jadi konsisten).
\end{solution}

\begin{solution}{exo:b2:fourier:2}
Kegenapannya membunuh $b_n$; lalu $a_0 = \frac{1}{\pi}\int_{-\pi}^\pi\abs
t\,\dd t = \pi$, dan untuk $n \geq 1$:
\[
a_n = \frac{2}{\pi}\int_0^\pi t\cos nt\,\dd t
= \frac{2}{\pi}\cdot\frac{(-1)^n - 1}{n^2}
= \begin{cases} -\frac{4}{\pi n^2} & n \text{ ganjil},\\ 0 & n
\text{ genap}, \end{cases}
\]
(lewat satu kali pengintegralan parsial). Karena itu
\[
\abs t = \frac{\pi}{2} - \frac{4}{\pi}\sum_{k\geq0}
\frac{\cos\bigl((2k+1)t\bigr)}{(2k+1)^2} ,
\]
dengan kekonvergenan \emph{normal} (sebab $\sum (2k+1)^{-2} < \infty$) --- seperti
diramalkan \cref{thm:b2:fourier:parseval}~(1) bagi fungsi \omterm{def:b2:metric:continuity}{kontinu} yang
$C^1$ sepotong-sepotong ini. Di $t = 0$:
\[
0 = \frac\pi2 - \frac4\pi\sum_{k\geq0}\frac{1}{(2k+1)^2}
\quad\Longrightarrow\quad
\sum_{k\geq0}\frac{1}{(2k+1)^2} = \frac{\pi^2}{8} .
\]
Memecah $\sum \frac{1}{n^2}$ menjadi bagian ganjil dan genap memberi $S =
\frac{\pi^2}{8} + \frac S4$, jadi $S = \frac{\pi^2}{6}$: yakni Basel lagi.
\end{solution}

\begin{solution}{exo:b2:fourier:3}
Kegenapannya: $b_n = 0$; lalu $a_0 = \frac{1}{\pi}\int_{-\pi}^{\pi} t^2 =
\frac{2\pi^2}{3}$; dan dua kali pengintegralan parsial memberi $a_n =
\frac{4(-1)^n}{n^2}$ (untuk $n \geq 1$). Karena itu
\[
t^2 = \frac{\pi^2}{3} + 4\sum_{n\geq1}
\frac{(-1)^n}{n^2}\cos nt
\qquad (\abs t \leq \pi),
\]
yang konvergen normal. Di $t = 0$: $0 = \frac{\pi^2}{3} +
4\sum\frac{(-1)^n}{n^2}$, yakni $\sum_{n\geq1}
\frac{(-1)^{n+1}}{n^2} = \frac{\pi^2}{12}$. Lalu Parseval:
\[
\frac{1}{2\pi}\int_{-\pi}^{\pi} t^4\,\dd t = \frac{\pi^4}{5}
= \frac{a_0^2}{4} + \frac12\sum_{n\geq1} a_n^2
= \frac{\pi^4}{9} + 8\sum_{n\geq1}\frac{1}{n^4} ,
\]
jadi $\sum \frac{1}{n^4} = \frac18\bigl(\frac{\pi^4}{5} -
\frac{\pi^4}{9}\bigr) = \frac{\pi^4}{90}$.
\end{solution}

\begin{solution}{exo:b2:fourier:4}
Bila lebih jauh $f$ bersifat $C^1$ sepotong-sepotong: maka Parseval (yang terbukti) memberi
$\norm f_2^2 = \sum\abs{c_n}^2 = 0$, dan kepositifan tegas
integral $\abs f^2$ yang \omterm{def:b2:metric:continuity}{kontinu} memaksa $f = 0$.

Untuk $f$ yang sekadar \omterm{def:b2:metric:continuity}{kontinu}, terimalah Parseval umum: buktinya
satu baris yang sama. (Garis besar rute kepadatannya: hampiran
trigonometri bertipe Fejér/Weierstrass menunjukkan bahwa polinomial itu
bersifat padat dalam $\norm\cdot_2$ di antara fungsi berkala yang \omterm{def:b2:metric:continuity}{kontinu}; dan karena $f
\perp$ semuanya, maka $\norm f_2^2 = \langle f, f - P\rangle \leq
\norm f_2\norm{f - P}_2$ untuk penghampir $P$, sehingga memaksa $\norm f_2 =
0$.)
\end{solution}

\begin{solution}{exo:b2:fourier:5}
Kegenapannya: $b_n = 0$; lalu
\[
a_n = \frac{2}{\pi}\int_0^\pi \cos(\alpha t)\cos(nt)\,\dd t
= \frac{2}{\pi}\cdot
\frac{(-1)^n\,\alpha\sin(\pi\alpha)}{\alpha^2 - n^2}
\]
(lewat hasil-kali-ke-jumlah, lalu integralkan; dengan $a_0 =
\frac{2\sin(\pi\alpha)}{\pi\alpha}$). Dirichlet di $t = \pi$ (yakni
titik kekontinuan perluasan berkalanya, yang nilai sepihaknya
sepakat berkat kegenapannya):
\[
\cos(\pi\alpha)
= \frac{\sin(\pi\alpha)}{\pi\alpha}
+ \sum_{n\geq1} \frac{2\alpha\sin(\pi\alpha)}{\pi(\alpha^2 -
n^2)}\,(-1)^n\cos(n\pi)
= \frac{\sin(\pi\alpha)}{\pi}\Bigl(\frac{1}{\alpha} +
\sum_{n\geq1}\frac{2\alpha}{\alpha^2 - n^2}\Bigr),
\]
dengan memakai $(-1)^n\cos n\pi = 1$. Membaginya dengan $\sin(\pi\alpha)/\pi$
memberi uraian kotangennya.
\end{solution}

\begin{solution}{exo:b2:fourier:6}
Mengiterasikan $c_n(f') = \iu n\,c_n(f)$ (sebanyak $k$ kali, lewat pengintegralan
parsial menyeberangi keping $C^{k}$-nya dengan nilai perbatasan yang bersesuaian) memberi
$c_n(f^{(k)}) = (\iu n)^k c_n(f)$. Adapun koefisien $f^{(k)}$ yang
\omterm{def:b2:metric:continuity}{kontinu} sepotong-sepotong bersifat terbatas (bahkan $\to 0$ menurut Bessel):
\[
\abs{c_n(f)} = \frac{\abs{c_n(f^{(k)})}}{\abs n^k}
= O\bigl(\abs n^{-k}\bigr) .
\]
\end{solution}

\begin{solution}{exo:b2:fourier:7}
Parseval bagi $f$ dan bagi $f'$ (keduanya sah: sebab $f$ bersifat $C^1$, dan $f'$
\omterm{def:b2:metric:continuity}{kontinu} sepotong-sepotong --- bahkan \omterm{def:b2:metric:continuity}{kontinu}):
\[
\frac{1}{2\pi}\int \abs f^2 = \sum_{n\neq0} \abs{c_n}^2
\quad (c_0 = 0 \text{ menurut hipotesis rerata-nol}),
\qquad
\frac{1}{2\pi}\int \abs{f'}^2 = \sum_{n\neq0} n^2\abs{c_n}^2 .
\]
Suku demi suku, $n^2\abs{c_n}^2 \geq \abs{c_n}^2$ untuk $\abs n \geq 1$:
jadi ketaksamaannya menyusul. Sedangkan kesamaannya memaksa $(n^2 - 1)\abs{c_n}^2 = 0$
untuk setiap $n$, yakni $c_n = 0$ untuk $\abs n \geq 2$: jadi $f(t) =
c_1\eu^{\iu t} + c_{-1}\eu^{-\iu t} = a\cos t + b\sin t$ (dalam bentuk
realnya); dan sebaliknya $f$ yang demikian memberi kesamaannya.
\end{solution}

\begin{solution}{exo:b2:fourier:8}
Dari \cref{exo:b2:fourier:1}, $S_{2N-1}(x) =
\frac4\pi\sum_{k=0}^{N-1}\frac{\sin((2k+1)x)}{2k+1}$. Di $x_N =
\frac{\pi}{2N}$, dengan $u_k = (2k+1)x_N = \frac{(2k+1)\pi}{2N}$:
\[
S_{2N-1}(x_N) = \frac{4}{\pi}\sum_{k=0}^{N-1}\frac{\sin u_k}{2k+1}
= \frac{4}{\pi}\sum_{k=0}^{N-1}\frac{\sin u_k}{u_k}\cdot
\frac{u_k}{2k+1}
= \frac{2}{\pi}\sum_{k=0}^{N-1}\frac{\sin u_k}{u_k}\cdot
\frac{\pi}{N},
\]
karena $\frac{u_k}{2k+1} = \frac{\pi}{2N}$. Adapun titik $u_k$ adalah
titik tengah $N$ subinterval $\intcc{0}{\pi}$ yang berpanjang
$\frac{\pi}{N}$: jadi jumlahnya adalah jumlah Riemann titik tengah bagi
$u \mapsto \frac{\sin u}{u}$ yang \omterm{def:b2:metric:continuity}{kontinu}, sehingga konvergen ke
\[
\frac{2}{\pi}\int_0^\pi \frac{\sin u}{u}\,\dd u
\approx \frac{2}{\pi}\times 1.8519 \approx 1.179 .
\]
Jumlah parsialnya di dekat lompatan melonjak melewati nilai $1$ sebesar $\approx
18\%$ setengah lompatannya untuk selamanya: yakni gejala Gibbs, yang dikuantifikasi.
\end{solution}

\begin{solution}{exo:b2:fourier:9}
Dari $\cos 3t = 4\cos^3t - 3\cos t$:
\[
\cos^3 t = \frac{3\cos t + \cos 3t}{4},
\qquad
\sin^2t\,\cos t = \cos t - \cos^3 t
= \frac{\cos t - \cos 3t}{4} .
\]
Masing-masingnya polinomial trigonometri, sehingga sama dengan deret
Fouriernya sendiri (berkat ketunggalan koefisiennya: sebab dua uraian
akan berselisih sebuah polinomial trigonometri yang semua
koefisiennya nol). Untuk $\cos^3t$: $a_1 = \frac34$, $a_3 =
\frac14$, sedangkan semua $a_n$ lainnya dan semua $b_n$ nol; lalu $c_{\pm1} =
\frac38$ dan $c_{\pm3} = \frac18$. Untuk $\sin^2t\cos t$: $a_1 =
\frac14$, $a_3 = -\frac14$; lalu $c_{\pm1} = \frac18$ dan $c_{\pm3} =
-\frac18$.
\end{solution}

\begin{solution}{exo:b2:fourier:10}
Lewat perhitungan langsung:
\[
c_n = \frac{1}{2\pi}\int_{-\pi}^{\pi}\eu^{(a - \iu n)t}\dd t
= \frac{\eu^{(a-\iu n)\pi} - \eu^{-(a - \iu n)\pi}}
{2\pi(a - \iu n)}
= \frac{(-1)^n\sinh(a\pi)}{\pi(a - \iu n)} ,
\]
dengan memakai $\eu^{\pm\iu n\pi} = (-1)^n$. Di $t = \pi$ perluasan
berkalanya melompat dari $\eu^{a\pi}$ ke $\eu^{-a\pi}$; lalu Dirichlet
(dengan jumlah parsial yang simetris) memberi
\[
\cosh(a\pi) = \sum_{n\in\Z}(-1)^n c_n\,
= \frac{\sinh(a\pi)}{\pi}\Bigl(\frac1a +
\sum_{n\geq1}\Bigl(\frac{1}{a - \iu n} +
\frac{1}{a + \iu n}\Bigr)\Bigr)
= \frac{\sinh(a\pi)}{\pi}\Bigl(\frac1a +
\sum_{n\geq1}\frac{2a}{a^2 + n^2}\Bigr) ,
\]
sebab bagian imajiner suku berpasangannya saling meniadakan. Bagilah dengan
$\sinh(a\pi)$:
\[
\coth(\pi a) = \frac{1}{\pi a} +
\sum_{n\geq1}\frac{2a}{\pi(a^2 + n^2)} ,
\]
yakni kembaran hiperbolik \cref{exo:b2:fourier:5}.
\end{solution}

\begin{solution}{exo:b2:fourier:11}
Kekomutatifannya: sulihkan $s = x - t$ lalu pakai kekalaan
integrannya. Untuk $c_n(f * g)$, integran $(x, t) \mapsto
f(x-t)g(t)\eu^{-\iu nx}$ bersifat \omterm{def:b2:metric:continuity}{kontinu}; jadi kedua integral
teriterasinya sepakat (sebab keduanya, sebagai fungsi batas atas
peubah luarnya, lenyap di ujung kirinya dan punya turunan yang
sama --- dan kekontinuannya beserta
\cref{thm:b2:integration:continuity} membenarkan penurunan
integral teriterasinya). Karena itu
\[
c_n(f*g) = \frac{1}{2\pi}\int_{-\pi}^{\pi} g(t)\,\eu^{-\iu nt}
\Bigl(\frac{1}{2\pi}\int_{-\pi}^{\pi}
f(x-t)\,\eu^{-\iu n(x-t)}\dd x\Bigr)\dd t
= c_n(f)\,c_n(g),
\]
dengan integral dalamnya bernilai $c_n(f)$ untuk setiap $t$ (lewat penyulihan
dan kekalaannya). Akhirnya \cref{lem:b2:fourier:kernel} mengatakan
$S_N(f)(x) = \frac{1}{2\pi}\int f(x+u)D_N(u)\dd u$; lalu
penyulihan $u \mapsto -t$ beserta kegenapan $D_N$ mengubah ini
menjadi $(f * D_N)(x)$.
\end{solution}

\begin{solution}{exo:b2:fourier:12}
\begin{enumerate}
  \item Dengan $w = r\eu^{\iu t}$:
        \[
        P_r(t) = 1 + 2\,\Re\frac{w}{1 - w}
        = \Re\frac{1 + w}{1 - w}
        = \frac{1 - \abs w^2}{\abs{1 - w}^2}
        = \frac{1 - r^2}{1 - 2r\cos t + r^2} > 0 .
        \]
        Rerataannya $1$: sebab pengintegralan suku demi suku deret yang
        konvergen normal itu hanya menyisakan $n = 0$.
  \item Untuk $\delta \leq \abs t \leq \pi$: berlaku $\cos t \leq
        \cos\delta$, jadi $P_r(t) \leq \frac{1 - r^2}{1 -
        2r\cos\delta + r^2}$, yang penyebutnya menuju $2 -
        2\cos\delta > 0$ sedangkan pembilangnya menuju $0$:
        sehingga konvergen \omterm{def:b2:funcseq:def}{seragam} ke $0$ di luar sembarang persekitaran
        $0$.
  \item Pengintegralan suku demi suku, berkat kekonvergenan yang normal dalam $t$,
        memberi $(f * P_r)(x) = \sum_n r^{\abs
        n}c_n(f)\eu^{\iu nx}$. Adapun hujah identitas
        hampirannya: dengan rerata $1$ dan kepositifannya,
        \[
        \abs{(f*P_r)(x) - f(x)}
        \leq \frac{1}{2\pi}\int_{-\pi}^{\pi}
        \abs{f(x-t) - f(x)}\,P_r(t)\,\dd t ,
        \]
        pecahlah di $\abs t = \delta$: maka paling banyak $\varepsilon$
        (menurut Heine) ditambah $2\norm f_\infty\sup_{\delta\leq\abs
        t\leq\pi}P_r \to \varepsilon$: jadi konvergen \omterm{def:b2:funcseq:def}{seragam} ketika
        $r \to 1^-$. Inilah \omterm{thm:b2:series:abel}{penjumlahan Abel} bagi deret
        Fourier --- yakni kembaran Fourier bagi teori perbatasan
        pada bab deret pangkat.
\end{enumerate}
\end{solution}

\begin{solution}{pb:b2:fourier:1}
\textbf{1.} Merata-ratakan \cref{lem:b2:fourier:kernel} atas $n = 0,
\dots, N-1$ (berkat kelinearan integralnya) memberi $\sigma_N(f)(x) =
\frac{1}{2\pi}\int f(x+u)F_N(u)\dd u$; dan tiap $D_n$ bererata
$1$, jadi $F_N$ bererata $1$.

\textbf{2.} Dengan $\eu^{\iu u} - 1 = 2\iu\,\eu^{\iu
u/2}\sin\frac u2$:
\[
\sum_{n=0}^{N-1}\sin\Bigl(\Bigl(n + \frac12\Bigr)u\Bigr)
= \Im\Bigl[\eu^{\iu u/2}\,\frac{\eu^{\iu Nu} - 1}{\eu^{\iu u}
- 1}\Bigr]
= \Re\,\frac{1 - \eu^{\iu Nu}}{2\sin\frac u2}
= \frac{1 - \cos Nu}{2\sin\frac u2}
= \frac{\sin^2\frac{Nu}2}{\sin\frac u2} .
\]
Membaginya dengan $N\sin\frac u2$:
\[
F_N(u) = \frac1N\sum_{n=0}^{N-1}
\frac{\sin\bigl((n+\frac12)u\bigr)}{\sin\frac u2}
= \frac{1}{N}\,
\frac{\sin^2\frac{Nu}{2}}{\sin^2\frac u2} \geq 0 .
\]

\textbf{3.} Pada $\delta \leq \abs u \leq \pi$: berlaku $\sin^2\frac u2
\geq \sin^2\frac\delta2$ dan $\sin^2\frac{Nu}2 \leq 1$, jadi $F_N
\leq \frac{1}{N\sin^2(\delta/2)} \to 0$ secara \omterm{def:b2:funcseq:def}{seragam} di sana.

\textbf{4.} Berkat rerata satuannya, $\sigma_Nf(x) - f(x) =
\frac{1}{2\pi}\int\bigl(f(x+u) - f(x)\bigr)F_N(u)\dd u$. Diberikan
$\varepsilon$, Heine menghasilkan $\delta$ dengan $\abs{f(x+u) - f(x)}
\leq \varepsilon$ untuk $\abs u \leq \delta$, secara seragam dalam $x$.
Lalu, dengan memakai $F_N \geq 0$ beserta rerata satuannya,
\[
\abs{\sigma_Nf(x) - f(x)}
\leq \varepsilon +
2\norm f_\infty\cdot\frac{1}{2\pi}
\int_{\delta\leq\abs u\leq\pi}F_N
\leq \varepsilon + \frac{2\norm
f_\infty}{N\sin^2\frac\delta2}
\leq 2\varepsilon
\]
untuk $N$ yang besar, secara seragam dalam $x$: itulah teorema Fejér.

\textbf{5.} Kernel $F_N$ bersifat genap dengan rerata $1$: jadi tiap paruh
$\intcc{0}{\pi}$ dan $\intcc{-\pi}{0}$ mengusung rerata $\frac12$.
Maka
\[
\sigma_Nf(x) - \frac{f(x^+)+f(x^-)}{2}
= \frac{1}{2\pi}\int_0^\pi\bigl(f(x+u) -
f(x^+)\bigr)F_N\,\dd u
+ \frac{1}{2\pi}\int_{-\pi}^0\bigl(f(x+u) -
f(x^-)\bigr)F_N\,\dd u ;
\]
lalu pada tiap paruhnya, pecahlah di $\abs u = \delta$ tempat limit
sepihaknya dekat sejauh $\varepsilon$, dan biarkan pertanyaan 3 membunuh bagian
jauhnya: maka kedua integralnya menuju $0$. Adapun batasnya: sebab $F_N \geq 0$
memberi $\abs{\sigma_Nf(x)} \leq
\frac{1}{2\pi}\int\abs{f(x+u)}F_N \leq \norm f_\infty$.

\textbf{6.} Tiap $\sigma_N(f)$ merupakan polinomial trigonometri
(yakni rata-rata $S_n(f)$ dengan $n < N$), dan $\sigma_N(f) \to f$
secara \omterm{def:b2:funcseq:def}{seragam}: itulah teorema Weierstrass trigonometri.

\textbf{7.} Syarat $c_n(f) = 0$ untuk setiap $n$ membuat setiap $S_n(f) = 0$,
sehingga setiap $\sigma_N(f) = 0$; jadi menurut Fejér, $f = \lim\sigma_Nf =
0$. Menerapkan ini pada sebuah selisih: maka fungsi berkala yang
\omterm{def:b2:metric:continuity}{kontinu} ditentukan oleh \omterm{def:b2:fourier:coefficients}{koefisien Fouriernya}.

\textbf{8.} Jumlah $S_Nf$ adalah proyeksi ortogonal $f$ ke
$\mathcal T_N$ (\cref{prop:b2:fourier:bessel}), jadi ia meminimumkan
$\norm{f - P}_2$ atas $P \in \mathcal T_N$; dan karena $\sigma_Nf
\in \mathcal T_{N-1} \subseteq \mathcal T_N$:
\[
\norm{f - S_Nf}_2 \leq \norm{f - \sigma_Nf}_2
\leq \norm{f - \sigma_Nf}_\infty \to 0
\]
(dengan ketaksamaan tengahnya karena rerata $\abs\cdot^2$ paling
banyak sebesar kuadrat supnya). Lalu Pythagoras $\norm f_2^2 = \sum_{\abs
n\leq N}\abs{c_n}^2 + \norm{f - S_Nf}_2^2$ melewatkan
limitnya: jadi Parseval bagi setiap $f$ berkala yang \omterm{def:b2:metric:continuity}{kontinu}.

\textbf{9.} Misalkan $t_1, \dots, t_p$ menyatakan lompatan $f$ dalam satu
kala, dan $M = \norm f_\infty$. Untuk $\eta$ yang kecil, definisikan $g = f$
di luar interval $\intoo{t_j - \eta}{t_j + \eta}$ dan lewat
tali busur afin menyeberangi tiap interval itu: maka $g$ bersifat \omterm{def:b2:metric:continuity}{kontinu},
berkala, $\norm g_\infty \leq M$, dan
\[
\norm{f - g}_2^2 \leq \frac{1}{2\pi}\,p\cdot(2M)^2\cdot2\eta
\leq \varepsilon^2
\]
untuk $\eta$ yang kecil. Bessel membuat $S_N$ menjadi kontraksi bagi
$\norm\cdot_2$, jadi
\[
\norm{f - S_Nf}_2
\leq \norm{f - g}_2 + \norm{g - S_Ng}_2 + \norm{S_N(g -
f)}_2
\leq 2\varepsilon + \norm{g - S_Ng}_2 ,
\]
lalu pertanyaan 8 memberi $\limsup_N\norm{f - S_Nf}_2 \leq
2\varepsilon$ untuk setiap $\varepsilon$: jadi $\norm{f - S_Nf}_2 \to
0$, dan Pythagoras menghasilkan Parseval bagi setiap $f$ yang
\omterm{def:b2:metric:continuity}{kontinu} sepotong-sepotong: sehingga ``yang diterima tanpa bukti'' pada
\cref{thm:b2:fourier:parseval} kini menjadi teorema.

\textbf{10.} Gelombang perseginya punya $\norm f_\infty = 1$, jadi
pertanyaan 5 memberi $\abs{\sigma_N(f)} \leq 1$ di mana-mana dan untuk
setiap $N$ --- jadi tanpa lonjakan, selamanya --- sedangkan
\cref{exo:b2:fourier:8} menunjukkan $\sup_xS_{2N-1}(f)(x) \to
\approx 1.179$. Adapun satu sifat yang bertanggung jawab: $F_N \geq 0$, jadi
$\sigma_Nf(x)$ merupakan \emph{rata-rata} terbobot nilai $f$
sehingga tak pernah dapat meninggalkan $\intcc{\min f}{\max f}$; sedangkan $D_N$ mengambil
nilai negatif, jadi $S_N$ dapat meninggalkannya.

\textbf{11.} Berlaku $\abs{\sin N\theta} \leq N\abs{\sin\theta}$ secara
induktif (sebab $\abs{\sin(N{+}1)\theta} \leq
\abs{\sin N\theta}\abs{\cos\theta} +
\abs{\cos N\theta}\abs{\sin\theta} \leq
(N+1)\abs{\sin\theta}$): jadi dengan $\theta = \frac u2$,
\[
F_N(u) = \frac{\sin^2\frac{Nu}2}{N\sin^2\frac u2}
\leq \frac{N^2\sin^2\frac u2}{N\sin^2\frac u2} = N .
\]
Adapun kecekungan $\sin$ pada $\intcc{0}{\frac\pi2}$ memberi
$\sin\frac u2 \geq \frac{u}{\pi}$ untuk $0 \leq u \leq \pi$, jadi
$F_N(u) \leq \frac{1}{N(u/\pi)^2} = \frac{\pi^2}{Nu^2}$.

\textbf{12.} Berkat kegenapannya dan pemecahan di $\frac1N$:
\[
\frac{1}{2\pi}\int_{-\pi}^{\pi}\abs uF_N
= \frac1\pi\int_0^\pi uF_N
\leq \frac1\pi\Bigl(\int_0^{1/N}uN\,\dd u +
\int_{1/N}^{\pi}\frac{\pi^2}{Nu}\,\dd u\Bigr)
= \frac{1}{2\pi N} + \frac{\pi\ln(\pi N)}{N} .
\]
Untuk $N \geq 2$: berlaku $\ln(\pi N) \leq \bigl(1 +
\frac{\ln\pi}{\ln2}\bigr)\ln N \leq 2.66\ln N$ dan
$\frac{1}{2\pi N} \leq \frac{\ln N}{N}$, jadi momennya bernilai
$\leq \frac{9\ln N}{N}$.

\textbf{13.} Bila $L$ konstanta \omterm{def:b2:metric:continuity}{Lipschitz} bagi $f$:
\[
\abs{\sigma_Nf(x) - f(x)}
\leq \frac{1}{2\pi}\int\abs{f(x+u) - f(x)}F_N(u)\dd u
\leq L\cdot\frac{1}{2\pi}\int\abs uF_N
\leq \frac{9L\ln N}{N},
\]
secara seragam dalam $x$.

\textbf{14.} Nilai $c_0(e_1) = 0$ memberi $S_0(e_1) = 0$, sedangkan
$S_n(e_1) = e_1$ untuk $n \geq 1$: jadi $\sigma_N(e_1) =
\frac{N-1}{N}e_1$, sehingga $\norm{\sigma_Ne_1 - e_1}_\infty =
\frac1N$. Jadi bahkan untuk sinyal utuh yang berpita terbatas ini lajunya tetap
$\frac1N$: Fejér menjenuh, persis seperti operator Bernstein
menjenuh di $\frac1n$ (yakni Voronovskaya, pada soal akhir pekan bab
barisan fungsi).

\textbf{15.} Bila $f$ lenyap pada $\intoo{x-\delta}{x+\delta}$,
maka $\sigma_Nf(x) = \frac{1}{2\pi}\int_{\delta \leq \abs u
\leq \pi}f(x+u)F_N(u)\dd u$, yang bernilai mutlak paling banyak
$\frac{\norm f_\infty}{N\sin^2(\delta/2)} \to 0$: jadi rata-rata Cesàro
di $x$ hanya melihat $f$ di dekat $x$.

\textbf{16.} Diberikan $\intcc ab$ dan $\varepsilon$, pilihlah
fungsi \omterm{def:b2:metric:continuity}{kontinu} berkala-$1$ yang afin sepotong-sepotong, sebut $\varphi^\pm$, dengan
$\varphi^- \leq \mathbf 1_{\intcc ab} \leq \varphi^+$ dan
$\int_0^1(\varphi^+ - \varphi^-) \leq \varepsilon$ (yakni trapesium
yang lerengnya di atas interval berpanjang total $\varepsilon$).
Maka
\[
\limsup_N\frac{\#\{n \leq N : x_n \in \intcc ab\}}{N}
\leq \lim_N\frac1N\sum_{n\leq N}\varphi^+(x_n)
= \int_0^1\varphi^+ \leq b - a + \varepsilon ,
\]
dan secara simetris $\liminf \geq b - a - \varepsilon$: jadi
proporsinya menuju $b - a$.

\textbf{17.} Untuk $f = \eu^{2\iu\pi k\cdot}$ dengan $k \neq 0$,
hipotesisnya memberi limit $0 = \int_0^1f$; sedangkan untuk $k = 0$ kedua
ruasnya bernilai $1$; lalu kelinearannya menangani setiap polinomial
trigonometri. Untuk fungsi \omterm{def:b2:metric:continuity}{kontinu} berkala-$1$, sebut $f$, dan $\varepsilon >
0$, pertanyaan 6 (yang diangkut lewat $t = 2\pi x$) menyediakan sebuah
polinomial trigonometri $P$ dengan $\norm{f - P}_\infty \leq
\varepsilon$:
\[
\Bigl|\frac1N\sum_{n\leq N}f(x_n) - \int_0^1f\Bigr|
\leq 2\varepsilon +
\Bigl|\frac1N\sum_{n\leq N}P(x_n) - \int_0^1P\Bigr|
\longrightarrow 2\varepsilon .
\]
Bersama pertanyaan 16: jadi $(x_n)$ terdistribusi merata.

\textbf{18.} Untuk $k \neq 0$ dan $\alpha$ yang irasional, dengan $w =
\eu^{2\iu\pi k\alpha} \neq 1$:
\[
\Bigl|\sum_{n=1}^{N}w^n\Bigr|
= \Bigl|\frac{w(w^N - 1)}{w - 1}\Bigr|
\leq \frac{2}{\abs{1 - w}} ,
\]
yakni batas yang tak bergantung pada $N$; jadi membaginya dengan $N$ memberi kriteria
Weyl, lalu pertanyaan 17 merampungkannya: sehingga $(\{n\alpha\})$
terdistribusi merata.

\textbf{19.} Setiap subintervalnya menerima proporsi asimtotik yang
sama dengan panjangnya, khususnya titik yang tak berhingga banyaknya: jadi
$(\{n\alpha\})$ bersifat padat. Adapun pemerataan lebih kuat: sebab sebuah
barisan dapat padat sambil menghabiskan hampir seluruh waktunya di satu
sudut (kepadatan mengatakan ke mana barisannya \emph{pergi},
sedangkan pemerataan mengatakan \emph{seberapa sering}).

\textbf{20.} Bilangan $2^n$ berdigit terdepan $1$ bila dan hanya bila $10^m \leq 2^n <
2\cdot10^m$ untuk suatu $m$, yakni bila $\{n\log_{10}2\} \in
\intco{0}{\log_{10}2}$. Keirasionalannya: sebab $\log_{10}2 =
\frac pq$ akan memberi $2^q = 10^p = 2^p5^p$, yang mustahil untuk $p
\geq 1$ berkat ketunggalan pemfaktorannya. Lalu teorema Weyl (pertanyaan 18
dengan $\alpha = \log_{10}2$; dan intervalnya yang setengah \omterm{def:b2:metric:topology}{terbuka} dijepit
di antara interval tertutup yang panjangnya berdekatan) memberi proporsinya
$\log_{10}2 \approx 0.301$: jadi digit pertama $2^n$ menuruti
hukum Benford.

\textbf{21.} Karena $z$ bersifat $C^1$, deret Fouriernya konvergen
normal dengan jumlah $z$ (\cref{thm:b2:fourier:parseval}~(1)),
dan $c_n(z') = \iu n\,c_n$. Lalu karena $\abs{z'} = \frac{L}{2\pi}$
bersifat tetap,
\[
\int_0^{2\pi}\abs{z'}^2\dd t
= 2\pi\Bigl(\frac{L}{2\pi}\Bigr)^{\!2}
= \frac{L^2}{2\pi} ,
\]
dan Parseval yang diterapkan pada $z'$ yang \omterm{def:b2:metric:continuity}{kontinu} memberi
$\frac{1}{2\pi}\int_0^{2\pi}\abs{z'}^2 =
\sum_n\abs{\iu nc_n}^2$, yakni $\frac{L^2}{2\pi} = 2\pi\sum_n
n^2\abs{c_n}^2$.

\textbf{22.} Parseval terpolarisasi $\frac{1}{2\pi}\int\conj
fg = \sum\conj{c_n(f)}c_n(g)$ menyusul dari Parseval yang diterapkan pada
$f + g$ dan $f + \iu g$ (lewat \omterm{def:b2:quadratic:def}{kesamaan polarisasinya}), yang keduanya
\omterm{def:b2:metric:continuity}{kontinu}. Dengan $f = z$ dan $g = z'$:
\[
A = \frac12\,\Im\int_0^{2\pi}\conj z\,z'
= \pi\,\Im\sum_n\conj{c_n}(\iu n c_n)
= \pi\sum_n n\abs{c_n}^2 .
\]

\textbf{23.} Menggabungkan pertanyaan 21--22:
\[
L^2 - 4\pi A
= 4\pi^2\sum_n n^2\abs{c_n}^2 -
4\pi^2\sum_n n\abs{c_n}^2
= 4\pi^2\sum_{n\in\Z}(n^2 - n)\abs{c_n}^2 \geq 0 ,
\]
karena $n^2 - n = n(n-1) \geq 0$ untuk setiap bilangan bulat. Sedangkan kesamaannya
memaksa $c_n = 0$ untuk setiap $n \notin \{0, 1\}$: jadi $z(t) = c_0 +
c_1\eu^{\iu t}$, yakni lingkaran berpusat $c_0$ dan berjari-jari
$\abs{c_1} = \frac{L}{2\pi}$ (berkat kelajuan tetapnya). Itulah bukti Hurwitz
bagi ketaksamaan isoperimetrik: $A \leq \frac{L^2}{4\pi}$,
dan hanya lingkaran.

\textbf{24.} Lingkaran berjari-jari $R$: $L = 2\pi R$ dan $A = \pi
R^2$, jadi $L^2 = 4\pi^2R^2 = 4\pi A$: yakni kesamaannya. Bujur sangkar bersisi
$a$: $L^2 = 16a^2 > 4\pi a^2 = 4\pi A$ (sebab $16 > 4\pi \approx
12.57$). Untuk $n$ yang negatif, $n^2 - n = n(n - 1)$ adalah hasil kali
dua bilangan bulat negatif: jadi positif --- sehingga ragam yang melilit mundur
memakan luas dua kali. Adapun kelajuan tetapnya masuk pada pertanyaan 21,
yang mengubah $\int\abs{z'}^2$ menjadi $\frac{L^2}{2\pi}$; sedangkan untuk
kelajuan yang tak tetap, Cauchy--Schwarz memberi $\int\abs{z'}^2 \geq
\frac{(\int\abs{z'})^2}{2\pi} = \frac{L^2}{2\pi}$, jadi
ketaksamaannya bertahan, dengan lingkaran tetap satu-satunya kasus
kesamaannya.

\textbf{25.} (i) Segalanya mengalir dari $F_N \geq 0$ (bersama rerata
satuannya dan pemusatannya); sedangkan $D_N$ punya rerata satuan dan pemusatan
ayunannya tetapi tanpa kepositifan, dan Gibbs adalah harganya.
(ii) Keterjumlahan Cesàro deret Fouriernya mengakibatkan
keterjumlahan Abelnya dengan jumlah yang sama (yakni Frobenius, yang terbukti pada
soal akhir pekan bab deret pangkat) --- jadi rute kernel Poisson
pada \cref{exo:b2:fourier:12} persis merupakan metode Abel.
(iii) Teorema Weyl hanya menuntut hampiran \omterm{def:b2:funcseq:def}{seragam}
(pertanyaan 6); sedangkan ketaksamaan isoperimetriknya menuntut Parseval
itu sendiri (pertanyaan 8, 21--22). (iv) Jilid Tahun ke-3 membuktikan
kelengkapannya: eksponensialnya membentuk basis Hilbert bagi $L^2$,
Parseval menjadi isometri ruang Hilbert, dan teorema Fejér
menjadi pernyataan bahwa isometri itu dapat dihitung
lewat rata-rata yang positif.

\end{solution}
